% Generated from validated Article Draft Result. Do not hand-edit; revise the structured draft instead.
\section{生成から行動へ――Tool UseがAPI層を作った}
\noindent\textbf{2023年には、言語モデルが文章を返すだけでなく、外部機能を選び、構造化された呼び出しへ接続する設計が製品・APIの中心へ入った。Toolformerの研究的発想とOpenAIのPlugins、Function Calling、GPTs/Assistantsの展開を、tool selectionと実行主体を分けて整理する。}\autocite{src-150e07b331df,src-4de1cea0a896,src-6b76b870b4ae,src-97786632d087,src-e17d0bf10781}

\subsection{文章を返すモデルから、外部処理を選ぶモデルへ}

tool useでは、モデルが「何を呼ぶべきか」を決めることと、外部システムが実際に副作用を伴う処理を実行することを分ける必要がある。自然言語の生成能力が高くても、実行権限、引数検証、失敗時の扱いは別のsystem layerである。\autocite{src-150e07b331df,src-4de1cea0a896,src-6b76b870b4ae,src-97786632d087,src-e17d0bf10781}


実装を抽象化すると、model outputから外部actionへは少なくとも「候補toolの選択」「machine-readableな引数生成」「host側の検証」「実行」「結果の観測」「次の生成」という段階がある。この分解を持つと、モデルの誤りと実行基盤の誤り、権限設計の誤りを別々に扱える。tool useがsystem designとして重要なのは、この境界が明示的に設計できるためである。\autocite{src-150e07b331df,src-4de1cea0a896,src-6b76b870b4ae,src-97786632d087,src-e17d0bf10781}


\subsection{研究上のtool useと、配備されたAPIは同じではない}

Toolformerはモデルが外部ツール利用を学習する研究方向を示した。一方、PluginsやFunction Callingは開発者が接続先と構造を定義する製品・API面の仕組みである。両者は同じ「tool use」という語で呼べても、研究上の問いと配備上の責任分界は異なる。\autocite{src-150e07b331df,src-4de1cea0a896,src-6b76b870b4ae,src-97786632d087,src-e17d0bf10781}


研究側が「どのtoolをいつ使うか」を学習問題として扱うのに対し、product側では、開発者がschema、権限、呼び出し先を定義し、モデルを制御されたinterfaceへ閉じ込めることが重要になる。この差は、agentという語を使うときにも残る。自律性を高めることと、安全に外部機能へ接続することは別の設計目標である。\autocite{src-150e07b331df,src-4de1cea0a896,src-6b76b870b4ae,src-97786632d087,src-e17d0bf10781}


\subsection{structured outputは信頼性の境界になる}

Function Callingのような構造化インターフェースは、自由文から外部実行へ直接飛ぶのではなく、モデル出力を機械可読な境界へ落とし込む考え方を前面に出した。この境界は利便性だけでなく、検証、再試行、権限管理を実装する場所にもなる。\autocite{src-4de1cea0a896,src-6b76b870b4ae,src-97786632d087,src-e17d0bf10781}


構造化された呼び出しでも、引数が妥当とは限らない。そのためhost側は型だけでなく、許容範囲、対象resource、実行権限、confirmationの要否を判断する必要がある。モデルを「toolを呼べる主体」として扱うほど、validationやauthorizationをモデル外へ置く設計原則が重要になる。\autocite{src-4de1cea0a896,src-6b76b870b4ae,src-97786632d087,src-e17d0bf10781}


toolの戻り値を次のprompt contextへ戻す設計では、外部世界の観測と生成を反復できる。これは静的な質問応答から一歩進む一方、失敗した実行結果や不完全な観測を次の判断材料にする危険も生む。loopを長くするほど、stop condition、error handling、human confirmationをどこに置くかがsystem behaviorを左右する。\autocite{src-4de1cea0a896,src-6b76b870b4ae,src-97786632d087,src-e17d0bf10781}


\subsection{tool accessがapplication-facing layerへ広がる}

年末までにGPTsやAssistantsへ連なる製品面が現れたことで、tool accessは個別デモではなくapplication-facing layerとして認識されやすくなった。2023年の転換は「モデルが何でも自律実行できるようになった」ことではなく、生成と実行の間に明示的なinterfaceが必要だと広く実装されたことである。\autocite{src-4de1cea0a896,src-6b76b870b4ae,src-97786632d087,src-e17d0bf10781}


このinterfaceがあると、application側はモデルを万能なブラックボックスとして扱わず、「利用可能なtool集合」と「許可された引数」を契約として持てる。結果としてmodel providerのAPI設計とapplication architectureが接続され、tool schemaはprompt engineeringだけでなくsoftware interface designの一部になる。\autocite{src-4de1cea0a896,src-6b76b870b4ae,src-97786632d087,src-e17d0bf10781}


ここからagent的なsystemへ発展するとしても、2023年のEvidenceから無制限な自律実行を一般化するべきではない。むしろtool/action layerの成立は、model reasoning、external execution、permission、observationを分けて組み合わせる必要性を示した。後のagent設計を理解する土台は、まずこの責任分界にある。\autocite{src-150e07b331df,src-4de1cea0a896,src-6b76b870b4ae,src-97786632d087,src-e17d0bf10781}


\begin{claimboundary}
tool-useの信頼性や有効性に関する評価は、研究・プロジェクト・vendorの条件に依存する。本稿では機構の成立と公開形態を扱い、それらの性能主張を独立再現済みの事実とはみなさない。\autocite{src-4de1cea0a896,src-6b76b870b4ae,src-97786632d087,src-e17d0bf10781}
\end{claimboundary}
